
On
notera que nous nous écartons ici de la terminologie usuelle (cf. [1]), où on appelait @@M0@@-anneau ce que nous appelons pré-@@M1@@-anneau, et @@M2@@-anneaux « spéciaux » nos @@M3@@-anneaux; la raison en est que seuls les @@M4@@-anneaux « spéciaux » sont utilisables en pratique.

\textbf{Exemples 2.5.}
\begin{enumerate}[label=\alph*)]
\item Nous avons déterminé toutes les structures de pré-@@M5@@-anneau de @@M6@@~; il résulte de 2.4.1 que seule celle qui est définie par @@M7@@ est susceptible d'être une structure de @@M8@@-anneau. On a alors @@M9@@, d'où~:
@@M10@@
@@M11@@
Par suite c'est bien une structure de @@M12@@-anneau. Si on désigne par @@M13@@ le polynôme à coefficients rationnels @@M14@@, on a donc~:
@@M15@@

\item Pour tout anneau @@M16@@, le pré-@@M17@@-anneau @@M18@@ est un @@M19@@-anneau. D'abord, @@M20@@. On vérifie les autres relations par la commutativité des diagrammes~:
@@M21@@
